\section{Delineamento}
\label{sec:delineamento}
\subsection{As interfaces e as exigências de coordenação}
\label{sec:as-interfaces-e-as-exigencias-de-coordenacao}

O conjunto de interfaces é um \emph{benchmark} de notação, como os \emph{benchmarks}
de resultado qualitativo do projeto de linguagens
\cite[p. 35]{erdweg2015}, entre os quais está o 7GUIs
\cite[p. 8:20]{lu2021}.
\notamudou{\textbf{Saíram as seções do projeto} sobre a revisão do JavaScript, o processamento de listas e a coordenação de eventos. As dimensões ficam na seção \ref{sec:as-dimensoes}, e a análise de cada interface vem depois deste delineamento.}

Das sete interfaces, quatro adaptam casos do 7GUIs: o Contador, o
\emph{Counter}; a Reserva de Passagem, o \emph{Flight Booker}, com a viagem de
ônibus no lugar do voo; a Planilha, o \emph{Cells}; e o Desenhador de
Círculos, o \emph{Circle Drawer} \cite[p. 17, 25, 40 e 49]{kiss2014}.
O Login segue o painel de \textcite[p. 534]{berry2020}, e o Placar ao
Vivo adapta a placares de jogos as cotações que uma fonte manda aos
clientes, de \textcite[p. 66]{meijer2012}.
A Busca traz o cancelamento de uma requisição já começada, cuja
resposta atrasada nunca aparece, comportamento assíncrono que nenhum
dos casos do 7GUIs testa, como reconhece \textcite[p. 77]{kiss2014}.
Cada interface exige da coordenação o que segue, com o que isso dá ao
usuário e, entre parênteses, o nome das exigências:

\begin{description}
\item[{Contador:}] um número e um botão que soma um a ele, sem exigência de
coordenação própria, o mínimo de código que cada tecnologia pede para
uma interface;
\item[{Busca com Sugestões:}] um campo de destino que sugere cidades a cada
termo novo de duas letras ou mais e cancela a busca do termo anterior,
para o usuário nunca ver as sugestões de um termo que já mudou
(assincronia, ciclo de vida da requisição, lista derivada e
histórico);
\item[{Login com Quarentena:}] um painel de nome e senha em que um novo
clique em Entrar cancela a tentativa anterior, para o usuário que errou
a senha tentar de novo sem esperar a resposta dela, e a sessão que
expira sozinha, com um relógio que para quando ela se encerra
(assincronia, ciclo de vida do relógio e estado derivado);
\item[{Placar ao Vivo:}] um painel de placares que uma fonte envia sem o
usuário pedir, numa conexão aberta ao entrar na página, para ele seguir
os jogos, e fechada ao pausar e ao sair dela, para ele ler a tabela
parada e não gastar rede com o que não vê (assincronia de uma fonte de
fora, ciclo de vida da conexão, estado derivado, lista derivada e
histórico);
\item[{Reserva de Passagem:}] um formulário de passagem de ônibus com o tipo
de viagem e duas datas, em que a data de volta só se habilita em ida e
volta e se marca como inválida quando é anterior à ida, para o usuário
saber por que não pode reservar, e o envio que ignora um novo clique
em Reservar até a resposta, para a mesma reserva não sair duas vezes
(estado derivado das duas datas, assincronia e ciclo de vida do envio
e histórico);
\item[{Planilha:}] uma coluna de oito células com o orçamento do mês e a
barra de fórmula, em que a alteração de uma célula se propaga só às
que dependem dela, para o usuário ver o saldo na hora em que muda um
gasto (estado derivado por um grafo que o usuário escreve, estado
compartilhado entre a grade e a barra e lista derivada);
\item[{Desenhador de Círculos:}] um \emph{canvas} em que o clique desenha um
círculo, o diâmetro ajustado num diálogo que guarda o círculo escolhido
e o desfazer e o refazer de cada passo, para o usuário corrigir o
desenho sem recomeçar (histórico, lista derivada e estado derivado).
\end{description}
\subsection{As tecnologias e os passos}
\label{sec:as-tecnologias-e-os-passos}

O núcleo tem duas tecnologias, o Web Component e o React, que
implementam as sete interfaces na primeira e na segunda rodada.
O Web Component é a forma do próprio navegador de criar elementos
personalizados, sem biblioteca: uma classe que estende \texttt{HTMLElement} e
que o navegador chama \cite{mozilla2026}.
Nele, o tratador de cada evento atribui o valor novo ao estado e escreve
também a atualização de cada elemento da tela que depende dele, a
orquestração por atribuição e a propagação manual.
No React, a tela se monta em JSX, marcação parecida com \gls{html}
escrita dentro do JavaScript \cite{meta2026c}.
O tratador também atribui o valor novo ao estado, mas não escreve a
atualização da tela: quando o estado se altera, o React reexecuta o
componente inteiro \cites[][ p. 289:38]{lee2025}{solidjs2026}, a
propagação por re-renderização.

O passo do Web Component ao React muda a propagação, de manual a
re-renderização, e deixa a orquestração por atribuição nos dois.
O nome de cada modo diz como a alteração do estado chega à tela, e não
como a tela se monta, e o mesmo passo muda duas coisas que o
delineamento não controla.
Uma é a montagem da tela.
Nas implementações deste trabalho, o esqueleto da tela do Web Component
fica em \gls{html} num texto, e só as partes que o estado repete ou põe
e tira da tela se montam elemento a elemento, com \texttt{createElement}.
No React, a tela inteira se monta em JSX.
A outra é o modelo de componente, de uma classe a uma função.
A análise atribui cada diferença pelo sinal escrito na notação: se ele
está na coordenação, na montagem da tela ou na ligação entre as duas.

O Web Component entra no núcleo como o extremo imperativo controlado,
sem biblioteca, e não como o representante do que se usa na prática.
Na pesquisa do Stack Overflow de 2025, o React é a segunda tecnologia
web mais usada entre quem respondeu à questão, com 44,7\%, e o jQuery,
na mesma célula do Web Component, a terceira, com 23,4\%
\cite{stackoverflow2025}.
Os respondentes vêm sobretudo dos canais do próprio site, o que favorece
os usuários mais ativos \cite{stackoverflow2025}.
O jQuery ficou de fora porque não organiza a tela em componentes: o
passo dele ao React acrescentaria um componente onde não havia nenhum.

Na terceira rodada, na Busca, no Login e no Placar ao Vivo, o núcleo
anda também na orquestração, sem trocar de tecnologia.
O \texttt{useReducer} leva o React à orquestração por mensagem: cada evento
vira uma ação, e uma função pura, o \emph{reducer}, recebe o estado e a ação
e devolve o estado seguinte \cite{meta2026d}.
O RxJS leva o Web Component e o React à orquestração por acumulação, em
que o estado se define sobre o fluxo de eventos, e, no último passo,
leva os derivados do React aos combinadores.

Depois da análise, o entorno do núcleo situa as células vizinhas, com
cada tecnologia num passo a partir de uma implementação do núcleo
(Figura \ref{fig:grade-do-nucleo-e-do-entorno}).
Na primeira rodada, nas sete interfaces, o Lit parte do Web Component,
e o Solid, do React.
O Lit se constrói sobre o Web Component: o componente continua uma
classe que o navegador chama como elemento personalizado, mas a
alteração de uma propriedade reativa agenda um ciclo de atualização,
que chama de novo o \texttt{render()} e atualiza a tela pelo \emph{template}
\cites{google2026d}{google2026e}.
O passo muda a propagação, de manual a re-renderização, e, com ela, a
montagem das partes repetidas, que passa do \texttt{createElement} ao \texttt{map} no
\emph{template}.
Alterar um objeto ou uma lista sem trocar a referência não dispara a
atualização.
A documentação recomenda, na maioria das aplicações, tratar os dados
como imutáveis, mas admite, quando o dado serve a um componente só,
alterá-lo no lugar e pedir a atualização com \texttt{requestUpdate()}
\cite{google2026e}.
É o que fazem o Desenhador de Círculos e a Planilha, que alteram a
lista e os mapas no lugar, como no Web Component.
No Solid, a tela também se monta em JSX, mas o componente não se
reexecuta, e a alteração de um \emph{signal} atualiza só os trechos da tela
que dependem dele \cites[][ p. 289:38]{lee2025}{solidjs2026}: o passo
muda a propagação, de re-renderização a granular.

\begin{figure}[htbp]
\caption{\label{fig:grade-do-nucleo-e-do-entorno}O núcleo, da Figura \ref{fig:grade-do-nucleo}, e o entorno na grade dos modos de orquestração e de propagação, com uma seta por passo.}
\centering
\begin{tikzpicture}[grade]
  % As tecnologias, uma linha da matriz por fileira de caixas: sequencial;
  % atribuição em duas (empilhadas na célula); mensagem; acumulação em
  % duas. As colunas: manual, por re-renderização, granular, por
  % combinadores.
  \matrix[celulas] (g) {
    |[entorno, largo] (eff)| \nagradecom{Web Component}{com Effection} & & & \\
    & |[entorno] (lit)| \nagrade{Lit} & & \\
    |[nucleo, largo] (wc)| \nagrade{Web Component} &
      |[nucleo] (react)| \nagrade{React} &
      |[entorno, estreito] (solid)| \nagrade{Solid} & \\
    & |[nucleo] (usered)| \nagradecom{React}{com \texttt{useReducer}} & & \\
    |[nucleo, largo] (wcrx)| \nagradecom{Web Component}{com RxJS} &
      |[nucleo] (reactrx)| \nagradecom{React}{com RxJS} & &
      |[nucleo, maislargo] (reactder)| \nagradecom{React}{com derivados no RxJS} \\
    & & & |[entorno, maislargo] (cycle)| \nagrade{Cycle.js} \\
  };

  % As bordas das células, tiradas das caixas: o meio de cada vão entre
  % colunas e entre fileiras, e 0,25 cm além das caixas nas bordas.
  \coordinate (x0) at ([xshift=-0.25cm]wc.west);
  \path (wc.east) -- (react.west) coordinate[midway] (x1);
  \path (react.east) -- (solid.west) coordinate[midway] (x2);
  \path (solid.east) -- (solid.east -| reactder.west) coordinate[midway] (x3);
  \coordinate (x4) at ([xshift=0.25cm]reactder.east);
  \coordinate (y0) at ([yshift=0.3cm]eff.north);
  \coordinate (y1) at ([yshift=0.3cm]lit.north);
  \path (wc.south) -- (wc.south |- usered.north) coordinate[midway] (y2);
  \path (usered.south) -- (reactrx.north) coordinate[midway] (y3);
  \coordinate (y4) at ([yshift=-0.3cm]cycle.south);

  % Os tons dos cantos e as linhas, atrás das caixas.
  \begin{scope}[on background layer]
    \fill[cantoimp] (x0 |- y0) rectangle (x1 |- y2);
    \fill[cantodec] (x1 |- y2) rectangle (x4 |- y4);
    \draw[linha] (x0 |- y1) -- (x4 |- y1);
    \draw[linha] (x0 |- y3) -- (x4 |- y3);
    \draw[linha] (x2 |- y0) -- (x2 |- y4);
    \draw[linha] (x3 |- y0) -- (x3 |- y4);
    \draw[linha] (x0 |- y0) rectangle (x4 |- y4);
  \end{scope}

  % Propagação, no alto: o título, os lados e os modos.
  \path (x0 |- y0) -- (x4 |- y0) coordinate[midway] (topo);
  \path (x0 |- y0) -- (x1 |- y0) coordinate[midway] (topoimp);
  \path (x1 |- y0) -- (x4 |- y0) coordinate[midway] (topodec);
  \node[modoalto] at (wc |- y0) {manual};
  \node[modoalto] at (react |- y0) {por re-renderização};
  \node[modoalto] at (solid |- y0) {granular};
  \node[modoalto] at (reactder |- y0) {por combinadores};
  \node[ladoalto] at (topoimp) {imperativo};
  \node[ladoalto] at (topodec) {declarativo};
  \node[eixoalto] at (topo) {Propagação (do estado à tela)};

  % Orquestração, à esquerda: os modos, os lados e o título, girados.
  \path (x0 |- y0) -- (x0 |- y1) coordinate[midway] (fseq);
  \path (x0 |- y1) -- (x0 |- y2) coordinate[midway] (fatr);
  \path (x0 |- y2) -- (x0 |- y3) coordinate[midway] (fmsg);
  \path (x0 |- y3) -- (x0 |- y4) coordinate[midway] (facu);
  \node[modoesq] at (fseq) {sequencial};
  \node[modoesq] at (fatr) {por atribuição};
  \node[modoesq] at (fmsg) {por mensagem};
  \node[modoesq] (acu) at (facu) {por acumulação};
  % A linha de base do lado girado, logo à esquerda do modo mais largo.
  \coordinate (xl) at ([xshift=-0.15cm]acu.west);
  \path (xl |- y0) -- (xl |- y2) coordinate[midway] (ladoimp);
  \path (xl |- y2) -- (xl |- y4) coordinate[midway] (ladodec);
  \path (xl |- y0) -- (xl |- y4) coordinate[midway] (ladoeixo);
  \node[ladoesq] at (ladoimp) {imperativo};
  \node[ladoesq] at (ladodec) {declarativo};
  \node[eixoesq] at (ladoeixo) {Orquestração (do evento ao estado)};

  % As divisas entre o imperativo e o declarativo, que atravessam os
  % rótulos dos lados.
  \draw[divisa] (x1 |- y4) -- ([yshift=0.95cm]x1 |- y0);
  \draw[divisa] (x4 |- y2) -- ([xshift=-0.3cm]xl |- y2);

  % 1.ª e 2.ª rodadas.
  \draw[primeiras] (wc) -- (react);
  % 3.ª rodada.
  \draw[terceira] (wc) -- (wcrx);
  \draw[terceira] (react) -- (usered);
  \draw[terceira] (usered) -- (reactrx);
  \draw[terceira] (wcrx) -- (reactrx);
  \draw[terceira] (reactrx) -- (reactder);
  % Entorno.
  \draw[doentorno] (wc) -- (eff);
  \draw[doentorno] ([xshift=-0.15cm]wc.north east) -- ([xshift=0.15cm]lit.south west);
  \draw[doentorno] (react) -- (solid);
  \draw[doentorno] (reactder) -- (cycle);

  % Legenda, em duas linhas centradas sob a grade.
  \path (x0 |- y4) -- (x4 |- y4) coordinate[midway] (base);
  \matrix[matrix of nodes, anchor=north, column sep=3pt, inner sep=0pt,
      nodes={anchor=base west, font=\scriptsize}] (leg1)
      at ([yshift=-0.35cm]base) {
    |[amostra, fill=cornucleo]| {} & núcleo &[0.45cm]
    |[amostra, draw=black!75, line width=0.6pt]| {} & entorno &[0.45cm]
    |[amostra, fill=corcantoimp, draw=black!25]| {} & imperativo nos dois eixos &[0.45cm]
    |[amostra, fill=corcantodec, draw=black!25]| {} & declarativo nos dois eixos \\
  };
  \matrix[matrix of nodes, anchor=north, column sep=3pt, inner sep=0pt,
      nodes={anchor=base west, font=\scriptsize}] (leg2)
      at ([yshift=-0.25cm]leg1.south) {
    |[amostra]| {} & passo da 1.ª e da 2.ª rodadas &[0.45cm]
    |[amostra]| {} & passo da 3.ª rodada &[0.45cm]
    |[amostra]| {} & passo do entorno \\
  };
  \draw[primeiras] (leg2-1-1.west) -- (leg2-1-1.east);
  \draw[terceira] (leg2-1-3.west) -- (leg2-1-3.east);
  \draw[doentorno] (leg2-1-5.west) -- (leg2-1-5.east);
\end{tikzpicture}

\small Fonte: elaborado pelo autor.
\end{figure}

Na Busca, no Login e no Placar ao Vivo, o Web Component com Effection
parte do Web Component da segunda rodada e muda a orquestração, de
atribuição a sequencial.
Na orquestração sequencial, o código espera o próximo evento e segue
em ordem.
O Cycle.js parte do React com os derivados no RxJS, o último passo da
terceira rodada, e muda de uma vez a biblioteca de fluxos, o modelo de
componente, os efeitos e a montagem da tela.
A diferença nesse passo não se atribui a uma variável só, como nos
casos de \textcite[p. 23]{kiss2014}, em que o \emph{toolkit} pesou mais
que o paradigma.
Se o prazo não der para tudo, o entorno sai antes do resto, uma
tecnologia por vez, nesta ordem: o Lit, na mesma célula do React; o
Web Component com Effection; o Solid; e o Cycle.js.
\subsection{As rodadas}
\label{sec:as-rodadas}

A avaliação se faz em três rodadas, cada uma mais estreita que a
anterior, como as de \textcite[p. 3]{kiss2014}: a primeira responde à
P1, e a segunda e a terceira, à P1.1 e à P1.2.
Cada rodada lê um eixo: a primeira, a propagação; a segunda, o
crescimento do código; a terceira, a orquestração
(Figura \ref{fig:interfaces-por-rodada}).

\begin{figure}[htbp]
\caption{\label{fig:interfaces-por-rodada}As sete interfaces rodada a rodada: a primeira constrói as 14 implementações e a segunda acrescenta de uma a três mudanças por interface, 13 ao todo, no mesmo passo do Web Component ao React; a terceira refatora o código da segunda na Busca com Sugestões, no Login com Quarentena e no Placar ao Vivo, nos cinco passos da Figura \ref{fig:grade-do-nucleo-e-do-entorno}. A caixa tracejada dá a ordem do que sai se o prazo apertar.}
\centering
% A figura das rodadas, variante B revista para o par com a grade: as
% interfaces na ordem de 3.1.1 (o Contador; as três em que pesa a
% combinação de eventos; as três em que pesa o gerenciamento de estado),
% as caixas do alto iguais às da grade e a seta grossa na legenda.
% A ordem de saída no corte, numa caixa tracejada depois do nome.
\newcommand{\corte}[1]{\tikz[baseline=(c.base)]\node[corte] (c) {#1};}
\begin{tikzpicture}[grade,
    tabela/.style={matrix of nodes, nodes in empty cells, row sep=0pt,
      column sep=0pt, inner sep=0pt, outer sep=0pt,
      nodes={anchor=center, inner sep=0.06cm, minimum height=0.62cm,
        text width=0.4cm, align=center, font=\footnotesize,
        execute at begin node={\hyphenpenalty=10000 }},
      column 1/.style={nodes={text width=2.8cm, align=left}},
      column 2/.style={nodes={text width=1.9cm}},
      column 3/.style={nodes={text width=3.75cm, align=left}},
      column 5/.style={nodes={text width=6.1cm, align=left}},
      row 1/.style={nodes={font=\scriptsize}},
      row 4/.style={nodes={minimum height=1.0cm}},
      row 9/.style={nodes={font=\scriptsize, minimum height=0.85cm}}},
    % A marca é um círculo preto, e não a cor do núcleo, que diz a
    % tecnologia.
    marca/.style={circle, fill=black, minimum size=5.5pt, inner sep=0pt},
    corte/.style={rounded corners=1pt, draw=black, line width=0.5pt,
      dash pattern=on 1.2pt off 0.8pt, inner sep=1pt, minimum size=0pt,
      text width=, font=\scriptsize, anchor=base},
    passos/.style={draw=black!40, rounded corners=2pt, inner sep=0pt},
    regua/.style={draw=black!15},
    forte/.style={draw=black!60},
    divisao/.style={draw=black!40},
    % Pela linha de base da última linha, para os títulos alinharem com
    % ou sem descendente.
    rodada/.style={anchor=base, yshift=0.3cm, inner sep=0pt, align=center,
      font=\footnotesize},
  ]
  \matrix[tabela] (m) {
    & implementa & mudanças & & os cinco passos, em cada interface \\
    Contador & & lista de contadores & & \\
    Busca com Sugestões & & espera de 300 ms\par \emph{combobox} pelo teclado & & \\
    Login com Quarentena & & quarentena\par \enquote{lembrar meu nome}\,\corte{2.º} & & \\
    Placar ao Vivo & & reconexão\par bolão\,\corte{1.º}\par frequência de 1, 5 ou 10\,s & & \\
    Reserva de Passagem & & volta empurrada pela ida & & \\
    Planilha & & referência circular & & \\
    Desenhador de Círculos & & arraste dos círculos\par salvar no servidor\par campo do diâmetro & & \\
    ao todo & 14 implementações & 13 mudanças, 26 \emph{commits} & & 15 refatorações \\
  };
  % Os títulos das rodadas sobre as colunas de cada uma.
  \path (m-1-2.north west) -- (m-1-2.north east)
    node[midway, rodada] (r1) {\textbf{1.ª rodada}\\\scriptsize propagação · P1};
  \path (m-1-3.north west) -- (m-1-3.north east)
    node[midway, rodada] {\textbf{2.ª rodada}\\\scriptsize crescimento · P1.1};
  \path (m-1-4.north west) -- (m-1-5.north east)
    node[midway, rodada] {\textbf{3.ª rodada}\\\scriptsize orquestração · P1.2};
  % O passo das duas primeiras rodadas, no alto delas, com as caixas e a
  % seta da grade.
  \path (m-1-2.west) -- (m-1-3.east) coordinate[midway] (meio);
  \node[nucleo, largo, anchor=south east, xshift=-0.25cm] (wc)
    at ([yshift=0.35cm]meio |- r1.north) {\nagrade{Web Component}};
  \node[nucleo, anchor=south west, xshift=0.25cm] (react)
    at ([yshift=0.35cm]meio |- r1.north) {\nagrade{React}};
  \draw[primeiras] (wc) -- (react);
  % Réguas: forte sob os títulos e sobre o total, leve entre as
  % interfaces, na base da célula mais alta da linha, a das mudanças.
  \draw[forte] (m-1-1.west |- m-1-3.south) -- (m-1-5.east |- m-1-3.south);
  \foreach \r in {2,...,7}
    \draw[regua] (m-\r-1.west |- m-\r-3.south) -- (m-\r-5.east |- m-\r-3.south);
  \draw[forte] (m-8-1.west |- m-8-3.south) -- (m-8-5.east |- m-8-3.south);
  \foreach \c in {2,3,4}
    \draw[divisao] (m-1-\c.west |- m-1-3.north) -- (m-9-\c.west |- m-9-3.south);
  % As marcas.
  \foreach \r in {2,...,8} \node[marca] at (m-\r-2) {};
  \foreach \r in {3,5} \node[marca] at (m-\r-4) {};
  \node[corte, anchor=center] at (m-4-4) {3.º};
  % Os cinco passos, num bloco só, porque valem iguais nas três.
  \path ([yshift=-0.08cm]m-3-5.west |- m-3-3.north) coordinate (pa)
        ([yshift=0.08cm]m-5-5.east |- m-5-3.south) coordinate (pb);
  \draw[passos, fill=white] ([xshift=0.04cm]pa) rectangle ([xshift=-0.02cm]pb);
  \path (pa) -- (pb) coordinate[midway] (pm);
  \node[anchor=center, inner sep=0pt, text width=6.0cm, align=left,
      font=\scriptsize, execute at begin node={\hyphenpenalty=10000 }]
      (lista) at (pm) {%
    \textit{orquestração}\par
    Web Component → Web Component com RxJS\par
    React → React com \texttt{useReducer}\par
    React com \texttt{useReducer} → React com RxJS\par\smallskip
    \textit{propagação}\par
    Web Component com RxJS → React com RxJS\par
    React com RxJS → React com derivados no RxJS};

  % Legenda.
  \matrix[matrix of nodes, anchor=north west, column sep=3pt, inner sep=0pt,
      nodes={anchor=base west, font=\scriptsize}] (leg)
      at ([yshift=-0.35cm]m-9-1.south west) {
    |[marca, yshift=0.5ex]| {} & a interface entra na rodada &[0.45cm]
    |[corte]| {1.º} & sai primeiro se o prazo apertar &[0.45cm]
    |[amostra]| {} & passo da 1.ª e da 2.ª rodadas \\
  };
  \draw[primeiras] (leg-1-5.west) -- (leg-1-5.east);
\end{tikzpicture}

\small Fonte: elaborado pelo autor.
\end{figure}

Na primeira rodada, o Web Component e o React implementam as sete
interfaces, 14 implementações, cada uma só com o que a tecnologia traz.
A orquestração é por atribuição nas duas, e a diferença entre elas está
na propagação, manual numa e por re-renderização na outra.

A segunda rodada muda o comportamento e faz o código crescer, porque é
nos problemas grandes que os efeitos da viscosidade mais aparecem
\cite[p. 161]{green1996d}.
Cada interface recebe funcionalidades, que mudam o que o usuário
percebe ou pode fazer, e requisitos de combinação de eventos, que deixam
a tela igual e mudam só como os eventos se combinam.
São 13 mudanças, numa ordem fixada antes da análise, cada uma num
\emph{commit} no Web Component e depois num no React, 26 \emph{commits} ao todo:

\begin{description}
\item[{Contador:}] a lista de contadores, com a remoção de um deles e o
total derivado, para o usuário contar várias coisas ao mesmo tempo;
\item[{Busca com Sugestões:}] a espera de 300 ms depois da última tecla,
para as sugestões não mudarem a cada letra, e o \emph{combobox} operável
pelo teclado \cite{w3c2026}, para quem não usa o mouse escolher uma
sugestão;
\item[{Login com Quarentena:}] a quarentena depois de três falhas seguidas,
que \textcite[p. 539]{berry2020} acrescentam ao painel na segunda
versão dele, para proteger a conta de quem tenta adivinhar a senha, e
o "lembrar meu nome", para o usuário não digitar o nome a cada
entrada;
\item[{Placar ao Vivo:}] a queda da conexão, que o navegador refaz sozinho
\cite[§ 9.2.1]{whatwg2026}, com o aviso "Reconectando…" até o
próximo lote, para o usuário saber que a tabela parou por queda, e não
por pausa; o bolão, com os pontos de cada palpite e o total
derivados, para o usuário ver quantos pontos fez; e a frequência de 1,
5 ou 10 s, que vai ao servidor numa requisição à parte, para o usuário
escolher de quanto em quanto tempo os placares mudam;
\item[{Reserva de Passagem:}] a data de volta empurrada ao dia seguinte à
ida quando a ida passa dela, para o usuário não ter de corrigir a
volta;
\item[{Planilha:}] a referência circular, que a Planilha acusa sem avaliar o
ciclo, a exigência com que \textcite[p. 54]{kiss2014} discute a
viscosidade do \emph{Cells}, para o usuário ver o erro em vez de a planilha
travar;
\item[{Desenhador de Círculos:}] arrastar os círculos, com o arraste no
desfazer, para o usuário mudar um círculo de lugar sem desfazer e
desenhar de novo; salvar o desenho num servidor simulado, na ordem
das mudanças, para o desenho não se perder; e o campo numérico do
diâmetro, ligado ao \emph{slider} nos dois sentidos, para o usuário
digitar o diâmetro exato.
\end{description}

Latências, tentativas, prazos e intervalos ficam fixos no domínio,
iguais nas duas tecnologias.

A terceira rodada muda só a forma: refatora o código da segunda rodada
da Busca, do Login e do Placar ao Vivo, as interfaces em que a
combinação de eventos pesa.
O comportamento não muda, e o teste de conformidade e as cenas da
segunda rodada servem de rede.
São cinco passos por interface, cada um num modo de um eixo, 15
refatorações ao todo:

\begin{itemize}
\item do Web Component ao Web Component com RxJS, a orquestração, de
atribuição a acumulação, com os derivados ainda escritos à mão no
\emph{listener};
\item do React ao React com \texttt{useReducer}, a orquestração, de atribuição a
mensagem;
\item do React com \texttt{useReducer} ao React com RxJS, a orquestração, de
mensagem a acumulação, com o mesmo \emph{reducer} no \texttt{scan};
\item do Web Component com RxJS ao React com RxJS, a propagação, de manual
a re-renderização;
\item do React com RxJS ao React com os derivados no RxJS, a propagação dos
derivados, do corpo do componente aos combinadores.
\end{itemize}

O Web Component, o React e os dois com RxJS formam um quadrado, em que
se lê se a diferença num eixo depende do modo no outro.
No React com RxJS, os derivados ficam no corpo do componente, forma
escolhida para que o passo seguinte mude só os derivados.

Se o prazo apertar, depois do entorno saem, nesta ordem, fixada antes
da análise: o bolão, o "lembrar meu nome" e o Login na terceira rodada.
Cada mudança sai inteira, nas duas tecnologias.
Não saem a primeira rodada nem a Busca e o Placar ao Vivo na terceira.
Também ficam a lista de contadores e o campo do diâmetro, as únicas
mudanças que trazem a remoção com a identidade dos que ficam e a
ligação nos dois sentidos.
\subsection{Os controles e o teste de conformidade}
\label{sec:os-controles-e-o-teste-de-conformidade}

Seis coisas poderiam diferenciar as implementações sem ser o modelo de
programação nem a notação, e as seis ficam fixas: a linguagem, a
especificação, as dependências, o domínio, a aparência e as versões.
Esse controle das variáveis é parte do planejamento da pesquisa, o
delineamento \cite[p. 54]{prodanov2013}.

Todas as implementações estão na web, em \emph{TypeScript} no modo estrito.
As de uma interface seguem a mesma especificação, com os mesmos textos,
regras e dados, e cada mudança da segunda rodada entra na especificação
antes do código.
Cada implementação usa só o que a tecnologia traz, sem bibliotecas além
das que fazem parte dela: o RxJS nos passos da terceira rodada e, no
entorno, o Effection e o Cycle.js.
As regras que não dependem da tela, como o formato das datas, as
mensagens, os servidores simulados e a análise das fórmulas da
Planilha, ficam num \emph{domínio} comum às implementações de cada
interface; o Contador não tem nenhuma.
A alteração do estado, como a propagação da Planilha e o desfazer do
Desenhador de Círculos, fica em cada implementação, para que a
diferença entre alterar no lugar e trocar por uma cópia não saia da
análise.

Na primeira rodada, duas regras valem para todas as interfaces.
As mensagens de situação ficam numa região que o leitor de tela anuncia
sem mover o foco \cite[critério 4.1.3]{w3c2024}, e a interface que sai da tela encerra o que abriu:
relógios, conexões, requisições e ouvintes.

Um \emph{teste de conformidade} por interface executa as mesmas verificações,
na mesma ordem, contra cada implementação, cada uma com o nome de um
cenário da especificação.
Ele se cola no console do navegador, com a implementação aberta, ou roda
por um comando que abre cada implementação no Chromium.
O teste confere também as mensagens de situação e, nas interfaces que
abrem relógios, conexões, requisições ou ouvintes, a limpeza ao sair da
tela.
\emph{Capturas} de tela dos estados escolhidos de cada interface, geradas
automaticamente, denunciam desvios de comportamento ou de aparência,
porque devem sair idênticas entre as tecnologias.
As implementações, os testes de conformidade e as capturas ficam num
repositório público, em \url{https://github.com/iquabius/unemat-tcc}.
Nele, as versões da linguagem, das bibliotecas e das ferramentas ficam
fixadas\footnote{No \texttt{package-lock.json}: TypeScript 6.0.3, React 19.3.0,
Solid 1.9.15 e Lit 3.3.3, com Vite 8.3.1 para o \emph{build} e Playwright
1.63.0 para as capturas e o teste de conformidade por comando.}, para que as implementações possam ser reproduzidas.

As implementações foram geradas com agentes de inteligência artificial,
em parte a partir de implementações que o autor escreveu antes em
JavaScript, com \emph{callbacks}, com classes e com as bibliotecas de fluxos
RxJS e xstream, e passam pelos mesmos controles: a especificação, o
teste de conformidade e as capturas.

Na aparência, que fica fixa num arquivo \gls{css} comum às
implementações de cada interface, a tela só cita as classes pelo nome.
As capturas de tela conferem que ela sai igual em todas as tecnologias.
Só o valor que depende do estado, como a posição de um menu, entra na
notação.

Com a aparência fora da notação, a avaliação não vê o estilo escrito
junto do componente, no próprio código que monta a tela.
Também não chega às dependências ocultas que
\textcite[p. 342]{green2006} acham no sistema formado pelo \gls{html}
e pelo \gls{css}.
\subsection{A avaliação}
\label{sec:a-avaliacao}

A avaliação cabe ao autor, que examina o código das implementações dimensão
a dimensão, sem participantes, pelas \glspl{dc}, apresentadas na seção
\ref{sec:dimensoes-cognitivas}.

Em cada rodada, a avaliação percorre as 14 \glspl{dc} da lista
completa \cite[p. 115-118]{blackwell2003}, definidas na seção
\ref{sec:as-dimensoes}, e não um subconjunto
escolhido antes, porque a escolha prévia pode atrapalhar, e o conjunto
inteiro traz resultados inesperados e úteis \cite[p. 332]{blackwell2001}.
Para a análise, o trabalho as agrupa pelo que descrevem.
Oito tratam da notação: \emph{nível de abstração}, \emph{proximidade de
descrição}, \emph{consistência}, \emph{concisão}, \emph{operações mentais difíceis},
\emph{dependências ocultas}, \emph{expressividade de papel} e \emph{notação secundária}.
Para \textcite[p. 334]{green2006}, duas delas, a expressividade de
papel e as operações mentais difíceis, sempre foram mal descritas.
Duas são do sistema, que o ambiente de edição agrava ou atenua: a
\emph{viscosidade} e a \emph{propensão a erros}.
As outras quatro pesam no ambiente de edição: o \emph{compromisso
prematuro}, a \emph{avaliação progressiva}, a \emph{visibilidade} e a
\emph{provisoriedade}.
Destas, a avaliação lê a face que aparece na notação, porque os próprios
autores suspeitam que há uma zona cinzenta entre as dimensões da notação e as
do ambiente \cite[p. 357]{green2006}.

Nenhuma dimensão é boa ou má em si, e a avaliação sempre leva em conta a
atividade a apoiar \cite[p. 115]{blackwell2003}.
Na primeira rodada, as atividades são a busca e a compreensão exploratória
do código pronto, e, na segunda e na terceira, a modificação \cite[p. 113]{blackwell2003}.

A análise é por interface, como os casos de
\textcite[p. 17-55]{kiss2014}, um por seção.
A seção de cada interface traz a especificação curta, as cenas, a
captura da tela e, logo abaixo, o módulo principal do Web Component e
do React, inteiro menos os \texttt{import}, com as linhas numeradas.
Seguem as 14 \glspl{dc}, cada uma em poucas linhas, que citam as linhas
da listagem em vez de repetir o trecho, com “sem diferença” quando os
modelos de programação não diferem nela.
Em cada \gls{dc}, a análise aponta se a diferença está na coordenação,
na montagem da tela ou na ligação entre as duas.
Na mesma seção ficam as rodadas da interface: a segunda, com o \emph{diff}
de cada mudança, e, na Busca, no Login e no Placar ao Vivo, a terceira,
com os cinco passos.
Cada passo entra como \emph{diff} quando ele é mais curto que a listagem
nova, senão como a listagem nova inteira, com as linhas que mudaram
realçadas, e o entorno segue a mesma regra depois da análise.
O anexo guarda só o domínio de cada interface.
Como as sete interfaces se leem inteiras, o avaliador não escolhe a
interface que dá a evidência de cada \gls{dc}, nem antes nem depois de
ver o código.

Na segunda e na terceira rodada, o \emph{diff} de cada mudança apoia a
viscosidade: a mesma mudança, preparada antes e feita em cada
implementação, como no teste de \textcite[p. 162]{green1996d}, mostra
quantas ações ela pede, e quantas outras restauram a consistência
\cite[p. 115]{blackwell2003}.
A concisão tem ainda o apoio de uma contagem das linhas não vazias e
sem comentários de cada implementação, sem o domínio e o estilo, e com
a montagem da tela onde ela estiver: no Web Component, o texto em
\gls{html} e os \texttt{createElement}; no React, o JSX.

A comparação fecha numa tabela-síntese, que reúne o perfil de cada
modelo de programação nas \glspl{dc} \cite[p. 107-108 e 118]{blackwell2003},
com uma linha por \gls{dc} e uma coluna por modelo: duas na primeira e
na segunda rodada, e seis com a terceira.
Cada célula dá a posição do modelo na \gls{dc} e de que interface vem a
evidência, e a mesma tabela se lê por exigência de coordenação.
Da primeira rodada para as outras duas, a linha de cada \gls{dc} diz se
a diferença entre os modelos de programação fica \emph{mantida},
\emph{ampliada}, \emph{reduzida}, \emph{anulada}, \emph{invertida} ou \emph{surgida}, e a célula
diz se a posição do modelo \emph{sobe}, \emph{desce} ou \emph{fica igual}.
Esses rótulos são deste trabalho, e não de uma fonte.

Quem avalia pelas \glspl{dc} o que construiu tem conflito de
interesse e, num sistema implementado ao longo de muito tempo, tende a não
ver os problemas que se acostumou a contornar \cite{dagit2006}.
O conflito diminui num sistema ainda em formação, e usar as \glspl{dc}
tarde no próprio sistema é melhor do que não usá-las \cite{dagit2006}.
\subsection{As dimensões}
\label{sec:as-dimensoes}
Cada uma das 14 \glspl{dc} vem abaixo com a definição curta e a
explicação de \textcite[p. 115-118]{blackwell2003}, na ordem do
agrupamento de “A avaliação”, com a página de cada uma.
Cada uma traz também um exemplo da web, escolhido por este trabalho,
como \textcite[p. 12-14]{kiss2014} faz com as dele.
Os exemplos ilustram a dimensão e não antecipam a posição de nenhuma
das tecnologias comparadas.

\begin{description}
\item[{Nível de abstração:}] os tipos e a disponibilidade de mecanismos de
abstração. Uma abstração muda a notação de base, e o sistema que
permite muitas é potencialmente difícil de aprender (p. 116-117).
Uma função dá nome a um cálculo que se repete, e quem a chama não
precisa saber como ele se faz.
\item[{Proximidade de descrição:}] o quanto a representação se aproxima do
domínio, ou o quanto a notação se relaciona com o resultado que ela
descreve (p. 117). Uma tabela em \gls{html}, escrita linha a linha
e célula a célula, tem a forma da tabela que aparece na tela.
\item[{Consistência:}] semânticas parecidas se expressam em formas
sintáticas parecidas. O usuário infere a estrutura pelos padrões da
notação, e a informação parecida que se apresenta de maneiras
diferentes fica obscurecida e compromete a usabilidade (p. 117). No \gls{dom}, um tratador de
clique se liga tanto pela propriedade \texttt{onclick} quanto por
\texttt{addEventListener}.
\item[{Concisão:}] a verbosidade da linguagem. Há notações prolixas, que
ocupam espaço demais na área de exibição (p. 117). A função de seta
escreve um \emph{callback} com menos símbolos que a expressão \texttt{function}.
\item[{Operações mentais difíceis:}] a demanda alta sobre os recursos
cognitivos, por exigir demais da memória de trabalho ou por pedir
objetivos aninhados em muitos níveis (p. 117). \emph{Callbacks} aninhados
em vários níveis obrigam o leitor a guardar de cabeça em que ponto de
cada um ele está.
\item[{Dependências ocultas:}] vínculos importantes entre entidades que não
se veem. Se uma entidade cita outra, que cita uma terceira, alterar a
terceira pode ter repercussões inesperadas, como nas células de uma
planilha e nos \emph{links} do \gls{html} (p. 116). Uma regra global de
\gls{css}, escrita longe do elemento, muda a aparência dele sem que
o código que monta o elemento a cite.
\item[{Expressividade de papel:}] o propósito de uma entidade se infere com
facilidade. A notação expressiva deixa ver por que o autor montou a
estrutura de certo jeito; nas outras, as entidades se parecem, e as
relações entre elas custam a aparecer (p. 116). Um \texttt{<button>} diz o
papel que um \texttt{<div>} com tratador de clique, sem o atributo \texttt{role},
esconde.
\item[{Notação secundária:}] informação a mais por meios que não a sintaxe
formal, como os comentários de uma linguagem de programação e as cores
e os formatos que acrescentam informação ao texto (p. 117). A
indentação e as linhas em branco dizem ao leitor o que o compilador
ignora.
\item[{Viscosidade:}] a resistência à mudança. Um sistema viscoso pede
muitas ações do usuário para um objetivo só, seja pela repetição da
mesma ação, seja pelas ações que restauram a consistência depois de
uma mudança (p. 115). Sem uma ferramenta de refatoração, renomear uma
propriedade do estado pede uma edição em cada lugar que a lê.
\item[{Propensão a erros:}] a notação convida ao erro, e o sistema protege
pouco; mecanismos de prevenção, como a declaração obrigatória de
identificadores, reduzem o problema (p. 116). O ouvinte que a
interface acrescenta à janela e esquece de remover ao sair da tela
continua a responder a eventos.
\item[{Compromisso prematuro:}] restrições na ordem de fazer as coisas,
como ter de declarar identificadores cedo demais (p. 116). No
\emph{TypeScript} estrito, o parâmetro de uma função declarada à parte
precisa do tipo escrito, mesmo antes de o programador saber tudo o que
ela vai receber.
\item[{Avaliação progressiva:}] o trabalho feito até ali pode ser conferido
a qualquer momento, como nos ambientes interpretados, em que o usuário
experimenta versões incompletas do programa (p. 118). O servidor de
desenvolvimento que recarrega a página a cada arquivo salvo mostra a
tela de uma versão ainda incompleta.
\item[{Visibilidade:}] a capacidade de ver os componentes com facilidade.
Sistemas que enterram a informação em encapsulamentos reduzem a
visibilidade (p. 116). O estado que uma biblioteca guarda dentro dela,
e que só se lê pelas funções dela, não aparece no código de quem a
usa.
\item[{Provisoriedade:}] o grau de compromisso com as ações ou as marcas.
Mesmo com restrições rígidas na ordem de fazer as coisas, ajuda poder
fazer ações provisórias, como registrar opções de projeto, esboçar ou
testar hipóteses do tipo “e se” (p. 117-118). Comentar um
trecho para provar outra versão é uma ação provisória no código.
\end{description}
